\section{Dagger Segal models and actual change of index}

Generic diagram-model results are distinguished from their universe-zero applications to the actual reversal and free-dagger-simplex indices. Strict mapping-space tests and projective derived fibrancy are kept separate.

\begin{resultblock}
\begin{proposition}[The actual free dagger simplex index]
\label{catalogue:B1}
The category of Part I free dagger simplices retains every enriched dagger functor between its objects.
It is equivalent to the explicit finite-path category of free dagger simplices.
The actual index functor $u$ from reversal simplices is bijective on objects and faithful;
on each reversal simplex map, its enriched dagger functor is the path-substitution map conjugated by the canonical comparison isomorphisms.
These statements identify the index used by the change-of-index adjunction with the manuscript's actual free-object index, rather than a subcategory having only monotone maps.
\end{proposition}

\Source{DaggerModels.PartII.DaggerSimplexPartIIndex.pathEquivalence;
DaggerModels.PartII.DaggerSimplexPartIIndex.actualU_map_conjugation;
DaggerModels.PartII.DaggerSimplexPartIIndex.actualU_obj_bijective;
DaggerModels.PartII.DaggerSimplexPartIIndex.actualU_faithful}{lean/DaggerModelsPartII/DaggerSimplexPartIIndex.lean}
\CatalogueNote{Scope}{The actual index and its comparison are constructed at the Part I free-simplex universe 0.}
\CatalogueNote{Manuscript context}{h1.lem.adjoints}
\CatalogueNote{Construction and reuse}{Native Part II index construction; the Part I free dagger simplex objects are imported through the pinned dependency.}
\end{resultblock}

\begin{resultblock}
\begin{proposition}[The projective Kan model on diagram categories]
\label{catalogue:B2}
For $C:\mathrm{Type}_u$ with a category structure in universe $u$, the category $\Fun(C,\sSet)$ has a projective model structure.
Fibrations and weak equivalences are detected at every object of $C$.
The generating cofibrations and generating trivial cofibrations arise by applying the evaluation left adjoints to simplicial boundaries and horns, respectively;
the corresponding small-object arguments are established.
Projective cofibrations are objectwise monomorphisms.
The model is combinatorial, and the arrow category of objectwise weak equivalences is accessible with a Smith solution set for every small family of arrows.
\end{proposition}

\Source{DaggerModels.PartII.ProjectiveKanDiagram.model_fibrations;
DaggerModels.PartII.ProjectiveKanDiagram.model_weakEquivalences;
DaggerModels.PartII.ProjectiveKanDiagram.boundaryGenerators_smallObjectArgument;
DaggerModels.PartII.ProjectiveKanDiagram.hornGenerators_smallObjectArgument;
DaggerModels.PartII.ProjectiveKanDiagram.combinatorial;
DaggerModels.PartII.ProjectiveKanDiagram.weakArrow_isAccessible;
DaggerModels.PartII.ProjectiveKanDiagram.weakSolutionSet}{lean/DaggerModelsPartII/ProjectiveKanDiagram.lean;
lean/DaggerModelsPartII/ProjectiveWeakAccessibility.lean}
\CatalogueNote{Scope}{Generic statement for C : Type u with Category.\{u\} C and SSet.\{u\}; Theorem B specializes to u = 0.}
\CatalogueNote{Manuscript context}{h1.lem.proj-rev, h1.lem.proj-dag}
\CatalogueNote{Construction and reuse}{Native Part II model construction, using the pinned Part I Kan model and Mathlib evaluation adjunctions.}
\end{resultblock}

\begin{resultblock}
\begin{proposition}[The injective Kan model and its generators]
\label{catalogue:B3}
For the same small category $C$, the category $\Fun(C,\sSet)$ has an injective model structure whose cofibrations are exactly diagram monomorphisms and whose weak equivalences are objectwise Kan weak equivalences.
A small family of quotient-subobject inclusions generates all monomorphisms as a weak-factorization left class.
The model is combinatorial and left proper, and every diagram is cofibrant.
For each fixed diagram, pointwise tensor and the literal mapping space form a Quillen adjunction;
maps into an injectively fibrant diagram give Kan mapping spaces.
\end{proposition}

\Source{DaggerModels.PartII.InjectiveKanDiagram.quotientGeneratedCofibrations_eq_cofibrations;
DaggerModels.PartII.InjectiveKanDiagram.model_cofibrations;
DaggerModels.PartII.InjectiveKanDiagram.model_weakEquivalences;
DaggerModels.PartII.InjectiveKanDiagram.combinatorial;
DaggerModels.PartII.InjectiveKanDiagram.model_isCofibrant;
DaggerModels.PartII.InjectiveKanDiagram.model_leftProper;
DaggerModels.PartII.InjectiveKanDiagram.model_tensorMapping_isQuillen}{lean/DaggerModelsPartII/InjectiveKanDiagramGeneration.lean;
lean/DaggerModelsPartII/InjectiveKanModel.lean}
\CatalogueNote{Scope}{Generic statement for C : Type u with Category.\{u\} C and SSet.\{u\}; Theorem B specializes to u = 0.}
\CatalogueNote{Manuscript context}{h1.lem.proj-rev, h1.lem.proj-dag}
\CatalogueNote{Construction and reuse}{Native Part II injective model construction; its quotient-subobject generation is proved rather than assumed.}
\end{resultblock}

\begin{resultblock}
\begin{proposition}[Literal Segal and walking-unitary tests]
\label{catalogue:B4}
For a reversal-indexed simplicial-set diagram $W$, the strict Segal condition requires every map induced by a spine inclusion in degree $n\geq2$ to be a Kan weak equivalence on literal mapping spaces.
The strict unitary completeness condition uses the selected zero vertex of the actual walking-unitary diagram.
Their conjunction is equivalent to testing every arrow of the manuscript's family $S_{\mathrm{rev}}$ by precomposition into $W$.
Injective fibrancy is a separate hypothesis when these strict tests are interpreted as the manuscript's unitarily complete dagger Segal space condition.
\end{proposition}

\Source{DaggerModels.PartII.SegalLocalizers.segalCondition;
DaggerModels.PartII.SegalLocalizers.unitaryCompletenessCondition;
DaggerModels.PartII.SegalLocalizers.unitarilyCompleteSegalCondition_iff_family}{lean/DaggerModelsPartII/SegalSpaceConditions.lean}
\CatalogueNote{Scope}{Strict test definitions and the family equivalence are stated for ReverseDaggerSpace.\{u\}; the localized Theorem B models use u = 0.}
\CatalogueNote{Manuscript context}{h1.not.S, h1.def.segal-complete}
\CatalogueNote{Construction and reuse}{Native Part II mapping-space formulation of the manuscript's spine and walking-unitary family.}
\end{resultblock}

\begin{resultblock}
\begin{proposition}[Strict and derived fibrancy of complete dagger Segal spaces]
\label{catalogue:B5}
In the injective localization at $S_{\mathrm{rev}}$, a diagram is fibrant exactly when it is injectively fibrant before localization and satisfies the literal Segal and unitary completeness tests.
In the projective localization, a diagram is fibrant exactly when it is objectwise Kan and admits an objectwise weak equivalence to an injectively fibrant diagram satisfying those strict tests.
For a projectively fibrant source, the same criterion may be checked on any chosen injectively fibrant replacement along an objectwise weak equivalence.
Thus the projective claim is a derived condition, not a strict test on an arbitrary diagram.
\end{proposition}

\Source{DaggerModels.PartII.SegalLocalizers.injectiveLocalObject_iff_strict;
DaggerModels.PartII.SegalLocalizers.projectiveLocalObject_iff_derived;
DaggerModels.PartII.SegalLocalizers.projectiveLocalObject_iff_strict_replacement;
DaggerModels.PartII.SegalLocalizers.injectiveSegalModel_isFibrant_iff;
DaggerModels.PartII.SegalLocalizers.projectiveSegalModel_isFibrant_iff}{lean/DaggerModelsPartII/InjectiveSegalLocality.lean;
lean/DaggerModelsPartII/SegalLocalizationModels.lean}
\CatalogueNote{Scope}{The strict/derived localization and fibrancy theorems are stated at SSet.\{0\}.}
\CatalogueNote{Manuscript context}{h1.lem.local-objects, h1.lem.rev-models}
\CatalogueNote{Construction and reuse}{Native Part II comparison of injective and projective derived mapping tests.}
\end{resultblock}

\begin{resultblock}
\begin{proposition}[Reversal localizations and the identity Quillen equivalence]
\label{catalogue:B6}
The projective and injective Kan models on reversal-indexed diagrams have the same objectwise weak equivalences, and every projective cofibration is an injective cofibration.
Localization at the manuscript's Segal and walking-unitary family, using functorial projective cofibrant resolutions on the projective side, yields left proper combinatorial model structures with the same localized weak equivalences.
Their identity adjunction is a Quillen equivalence.
The injective localization satisfies the literal pushout-product axiom for diagram cofibrations and cofibrations of simplicial sets, with a weak-equivalence conclusion when either input morphism is a weak equivalence.
\end{proposition}

\Source{DaggerModels.PartII.ProjectiveInjectiveComparison.cofibrations_le;
DaggerModels.PartII.ProjectiveInjectiveComparison.localEquivalences_eq;
DaggerModels.PartII.SegalLocalizers.segalIdentity_isQuillenEquivalence;
DaggerModels.PartII.SegalLocalizers.injectiveSegalModel_combinatorial;
DaggerModels.PartII.SegalLocalizers.injectiveSegalModel_leftProper;
DaggerModels.PartII.SegalLocalizers.injectiveSegalModel_pushoutProduct_sm7}{lean/DaggerModelsPartII/ProjectiveInjectiveComparison.lean;
lean/DaggerModelsPartII/SegalLocalizationModels.lean}
\CatalogueNote{Scope}{The underlying projective/injective comparison is generic in u; the reversal Segal localization is at SSet.\{0\}.}
\CatalogueNote{Manuscript context}{h1.lem.rev-models, h1.lem.localized-simplicial}
\CatalogueNote{Construction and reuse}{Native Part II localizations; projective resolution and literal localizer weak classes are identified before the Quillen comparison.}
\end{resultblock}

\begin{resultblock}
\begin{proposition}[Localization on the actual free dagger simplex index]
\label{catalogue:B7}
The actual-index localizing family consists of the literal left Kan images of the chosen functorial projective resolutions of $S_{\mathrm{rev}}$.
Its projective localization is left proper and combinatorial.
An actual-index diagram is fibrant there precisely when every value is Kan and its restriction along $u$ satisfies derived unitary completeness and Segal conditions.
Localizing the injective actual-index model at the same family gives a left proper combinatorial model with the same localized weak equivalences;
the identity is a Quillen equivalence.
The injective localization also satisfies the actual diagram pushout-product axiom.
\end{proposition}

\Source{DaggerModels.PartII.ActualIndexProjectiveModel.actualSegalModel_isFibrant_iff;
DaggerModels.PartII.ActualIndexProjectiveModel.actualSegalModel_combinatorial;
DaggerModels.PartII.ActualIndexProjectiveModel.actualSegalModel_leftProper;
DaggerModels.PartII.ActualIndexProjectiveModel.actualSegalModels_weakEquivalences_eq;
DaggerModels.PartII.ActualIndexProjectiveModel.actualSegalIdentity_isQuillenEquivalence;
DaggerModels.PartII.ActualIndexProjectiveModel.actualInjectiveSegalModel_pushoutProduct_sm7}{lean/DaggerModelsPartII/ActualSegalLocalization.lean;
lean/DaggerModelsPartII/ActualInjectiveSegalLocalization.lean}
\CatalogueNote{Scope}{The actual free dagger simplex localization is constructed at SSet.\{0\}.}
\CatalogueNote{Manuscript context}{h1.def.model, h1.prop.fibrant, h1.thm.projection}
\CatalogueNote{Construction and reuse}{Native Part II actual-index construction; its localizing maps are the named Kan images of resolved manuscript maps.}
\end{resultblock}

\begin{resultblock}
\begin{proposition}[Localized change of index and weak-equivalence detection]
\label{catalogue:B8}
For the actual functor $u$ from reversal simplices to Part I free dagger simplices, left Kan extension and restriction form a Quillen adjunction between the specified projective Segal localizations.
Restriction reflects localized weak equivalences between fibrant actual-index diagrams.
A subsequent theorem strengthens this to every arrow:
an actual-index map is a localized weak equivalence exactly when its restriction is one.
Combining this detection with the cofibrant-unit result proves that the same actual adjunction is a Quillen equivalence.
The assertion concerns these localizations and this fixed $u$, not arbitrary changes of index.
\end{proposition}

\Source{DaggerModels.PartII.ActualIndexProjectiveModel.actualLeftAdjunction_localized_isQuillen;
DaggerModels.PartII.ActualIndexProjectiveModel.actualRestriction_reflectsLocalizedWeakEquivalences;
DaggerModels.PartII.ActualLocalizedEquivalence.actualRestriction_weakEquivalence_iff;
DaggerModels.PartII.ActualLocalizedEquivalence.actualLeftAdjunction_isQuillenEquivalence}{lean/DaggerModelsPartII/ActualLocalizedQuillen.lean;
lean/DaggerModelsPartII/ActualLocalizedEquivalence.lean}
\CatalogueNote{Scope}{The actual localized adjunction and detection theorem are at SSet.\{0\}.}
\CatalogueNote{Manuscript context}{h1.thm.adjunction, h1.thm.conservative, h1.thm.main}
\CatalogueNote{Construction and reuse}{Native Part II comparison for the literal Part I free-object index functor.}
\end{resultblock}

\begin{resultblock}
\begin{proposition}[Representable units and all projectively cofibrant units]
\label{catalogue:B9}
For every reversal representable indexed by $m$, the unit of the actual left Kan extension--restriction adjunction is a trivial cofibration in the injective Segal localization.
The proof constructs a length filtration of the native free dagger nerve and establishes local triviality at each stage of length at least two.
The unit theorem then extends to every diagram cofibrant in either the original or localized projective model.
This is the specific cofibrant-input statement needed for the derived unit of the localized adjunction;
it does not assert the strict unit is a local equivalence for arbitrary input.
\end{proposition}

\Source{DaggerModels.PartII.ClassifyingTheorem.lengthStage_trivialCofibration;
DaggerModels.PartII.ClassifyingTheorem.actualRepresentableUnit_trivial;
DaggerModels.PartII.CofibrantUnit.actualUnit_trivial_of_projectiveCofibrant;
DaggerModels.PartII.CofibrantUnit.actualUnit_trivial_of_localizedProjectiveCofibrant}{lean/DaggerModelsPartII/ClassifyingTheorem.lean;
lean/DaggerModelsPartII/CofibrantUnit.lean}
\CatalogueNote{Scope}{The representable and cofibrant-unit results are at SSet.\{0\}.}
\CatalogueNote{Manuscript context}{h1.prop.classifying, h1.lem.unit-cofibrant, h1.prop.derived-unit}
\CatalogueNote{Construction and reuse}{Native Part II length-filtration and projective-cell arguments for the actual adjunction unit.}
\end{resultblock}

\begin{resultblock}
\begin{proposition}[The literal simplicial enrichment behind Theorem B]
\label{catalogue:B10}
The reversal diagram category carries a native simplicial-category structure whose enriched Hom bifunctor is naturally isomorphic to the literal mapping space with $k$-simplices given by natural transformations $A\times\Delta^k\to W$.
The comparison sends an ordinary morphism to its canonical zero simplex.
Pointwise tensor and cotensor satisfy parameterized tensor--mapping and tensor--cotensor adjunctions, the corresponding internal-Hom exponential isomorphisms, and the stipulated naturality in diagram and simplicial-set variables.
These are explicit structure fields constructed in Lean, beyond the seven model-theoretic clauses of the core Theorem B statement.
\end{proposition}

\Source{DaggerModels.PartII.DiagramSimplicialEnrichment.homComparison_ordinary;
DaggerModels.PartII.theoremBSimplicialStructure}{lean/DaggerModelsPartII/DiagramSimplicialEnrichment.lean;
lean/DaggerModelsPartII/ProofTheoremB.lean}
\CatalogueNote{Scope}{The native simplicial enrichment and full literal structure are constructed for ReverseDaggerSpace.\{0\} and SSet.\{0\}.}
\CatalogueNote{Manuscript context}{h1.not.discrete, h1.lem.proj-rev, h1.lem.localized-simplicial}
\CatalogueNote{Construction and reuse}{Native Part II literal-enrichment implementation; theoremBSimplicialStructure supplies all fields of FullTheoremBStatement.LiteralSimplicialStructure.}
\end{resultblock}
